\subsection{Stellar algebras}

DEFINITION 3. --- A stellar algebra is an involutive Banach algebra A such that \(\|x\|^{2}=\|x^{*}x\|\) for all \(x\in A\) .

If A and B are stellar algebras, a morphism, or morphism of stellar algebras, from A to B is a morphism of involutive algebras from A to B. An isomorphism from A to B is an isomorphism of involutive algebras from A to B.

Some authors speak of "C*-algebra".

Let A be a stellar algebra. A stellar subalgebra of A is a closed involutive subalgebra of A.

Examples. --- 1) The involutive Banach algebra of continuous endomorphisms of a complex Hilbert space (example 3 of I, p.~100) is a stellar algebra (EVT, V, p.~39, prop. 2).

\begin{enumerate}
\def\labelenumi{\arabic{enumi})}
\setcounter{enumi}{1}
\tightlist
\item
  Let X be a locally compact topological space. The involutive Banach algebra \(\mathcal{C}_{b}(X)\) of continuous bounded complex-valued functions on X (example 1 of I, p.~100) is a stellar algebra. Indeed, for every function \(f \in \mathcal{C}_b(\mathrm{X})\) , one has \(f^*f = |f|^2\) , and hence \(\| f^*f \| = |||f|^2 \| = \| f\|^2\) .
\end{enumerate}

Let \(\mathrm{A} = \mathcal{C}_0(\mathrm{X})\) be the involutive Banach subalgebra of functions tending to 0 at infinity. It is a stellar subalgebra of \(\mathcal{C}_b(\mathrm{X})\) . For every function \(f \in \mathrm{A}\) , one has \(\| f \| = \varrho(f)\) since \(\mathrm{Sp}_{\mathrm{A}}'(f) = f(\mathrm{X}) \cup \{0\}\) .

Let X and Y be locally compact topological spaces. For every proper partial map \(\varphi\) from X into Y (def. 1 of I, p.~33), the algebra morphism \(\varphi^{*}\) of \(\mathcal{C}_{0}(\mathrm{Y})\) into \(\mathcal{C}_{0}(\mathrm{X})\) (prop. 3 of I, p.~34) is a morphism of involutive algebras. Conversely, every morphism of stellar algebras \(\pi\colon\mathcal{C}_{0}(\mathrm{Y})\to\mathcal{C}_{0}(\mathrm{X})\) is of this form (loc. cit.).

\begin{enumerate}
\def\labelenumi{\arabic{enumi})}
\setcounter{enumi}{2}
\item
  Let X be a compact topological space and \(x_{0} \in X\) a fixed element of X. The subalgebra \(\mathcal{C}'(X)\) of \(\mathcal{C}(X)\) of continuous functions \(f: X \to C\) such that \(f(x_{0}) = 0\) is a stellar algebra.
\item
  Let X be a separated topological space and \(\mu\) a positive measure on X. The involutive Banach algebra \(\mathrm{L}^{\infty}(\mathrm{X},\mu)\) is a commutative unital stellar algebra.
\item
  Let \((\mathbf{A}_{i})\) be a family of stellar algebras. The product involutive Banach algebra A of the \(A_{i}\) (example 5 of I, p.~101) is a stellar algebra, called the product stellar algebra of the \(A_{i}\) .
\item
  Let A be a stellar algebra. If B is a closed involutive subalgebra of A, then B is a stellar algebra. It will be seen (V, to appear) that every stellar algebra is isomorphic to a closed involutive subalgebra of the stellar algebra of endomorphisms of a Hilbert space (example 1).
\item
  Let A be a stellar algebra. If M ⊂ A is any subset, then the closed involutive subalgebra of A generated by M is a stellar algebra, called the stellar subalgebra of A generated by M. If A is moreover unital, then the closed unital involutive subalgebra generated by M is a unital stellar algebra, called the unital stellar subalgebra of A generated by M.
\item
  In general, the involutive Banach algebra \(\mathcal{M}^{1}(G)\) (example 4 of I, p.~100) is not a stellar algebra (I, p.~181, exerc. 8).
\end{enumerate}

Lemma 6. --- Let A be a Banach algebra equipped with an involution satisfying

\[
\| x \| ^ {2} \leqslant \| x ^ {*} x \|\tag{1}
\]

for all \(x \in \mathbf{A}\) . Then A is a stellar algebra.

Let \(x \in \mathbf{A}\) . We then have \(\|x\|^2 \leqslant \|x^*\| \cdot \|x\|\) , whence \(\|x\| \leqslant \|x^*\|\) . By exchanging the roles of \(x\) and \(x^*\) we see that \(\|x\| = \|x^*\|\) . Thus \(\|x^*x\| \leqslant \|x^*\|\|x\| \leqslant \|x\|^2\) and the hypothesis implies the equality \(\|x\|^2 = \|x^*x\|\) .

Lemma 7. --- Let A be a stellar algebra.

\begin{enumerate}
\def\labelenumi{\alph{enumi})}
\tightlist
\item
  The regular representation \(\gamma\) of A (def. 1 of I, p.~16) is isometric, that is,
\end{enumerate}

\[
\| x \| = \sup _ {\| y \| \leqslant 1} \| x y \|,
\]

for all \(x\in \mathbf{A}\) .

\begin{enumerate}
\def\labelenumi{\alph{enumi})}
\setcounter{enumi}{1}
\tightlist
\item
  For every \(x \in A_{h}\) , one has
\end{enumerate}

\[
\varrho (x) = \| x \|.\tag{2}
\]

Let \(x \in \mathrm{A}\) . We have \(\sup_{\| y \| \leqslant 1} \| xy \| \leqslant \| x \|\) . To prove that \(\| x \| \leqslant \sup_{\| y \| \leqslant 1} \| xy \|\) , we may assume \(\| x \| = 1\) . Then the element \(y = x^*\) satisfies \(\| y \| = \| x^* \| = 1\) , and \(\| xy \| = \| x \|^2 = 1\) , whence the assertion \(a\) ).

Suppose that \(x\) is Hermitian. It follows \(\| x^2 \| = \| x^* x \| = \| x \|^2\) , whence by induction \(\| x^{2^n} \|^{2^{-n}} = \| x \|\) for every integer \(n \geqslant 1\) , whence the assertion \(b)\) by prop. 1 of I, p.~20.

Remarks. --- 1) Let A be a unital stellar algebra. Then we have

\[
\left\| 1 \right\| ^ {2} = \left\| 1 ^ {*} 1 \right\| = \left\| 1 \right\|,
\]

therefore the norm \(\|1\|\) is zero or equal to 1. If A ≠ \{0\}, we deduce \(\|1\| = 1\) . Consequently, for every unitary element u of A, we have \(\|u\| = \|u^{*}u\|^{1/2} = 1\) .

\begin{enumerate}
\def\labelenumi{\arabic{enumi})}
\setcounter{enumi}{1}
\tightlist
\item
  Let A be a normed involutive algebra. If \(\|x\|^{2}=\|x^{*}x\|\) for all \(x\in A\) , the completion \(\widehat{A}\) of A is a stellar algebra.
\end{enumerate}

PROPOSITION 2. --- Let A be an involutive Banach algebra, let B be a stellar algebra, and let π be a morphism of involutive algebras from A to B. Then one has \(\|\pi(x)\| \leqslant \|x\|\) for all \(x \in A\) , and in particular π is continuous.

For every \(x \in \mathrm{A}\) , one has \(\mathrm{Sp}_{\mathrm{B}}^{\prime}(\pi(x)) \subset \mathrm{Sp}_{\mathrm{A}}^{\prime}(x)\) , hence \(\varrho(\pi(x)) \leqslant \varrho(x) \leqslant \|x\|\) . Since \(\pi(x^{*}x) \in \mathrm{B}_{h}\) , one has \(\|\pi(x^{*}x)\| = \varrho(\pi(x^{*}x))\) (formula (2)), hence

\[
\| \pi (x) \| ^ {2} = \| \pi (x ^ {*} x) \| = \varrho (\pi (x ^ {*} x)) \leqslant \| x ^ {*} x \| = \| x \| ^ {2}.
\]

PROPOSITION 3. --- Let A be a stellar algebra and let \(\widetilde{\mathrm{A}}\) be the involutive algebra obtained from A by adjoining a unit element. There exists a unique norm on \(\widetilde{\mathrm{A}}\) extending that of A and making \(\widetilde{\mathrm{A}}\) a stellar algebra.

The uniqueness of such a norm follows from Prop. 2. We now show its existence. Let us denote by \(\widetilde{e}\) the identity element of \(\widetilde{\mathrm{A}}\) .

First, suppose that A has a unit element e. The product of the involutive normed algebras A and \(\mathbf{C}(\widetilde{e}-e)\) is identified with \(\widetilde{A}\) and is a stellar algebra (example 5). The norm on \(\widetilde{A}\) extends that of A, whence the assertion.

Now suppose that A has no identity element. For every \(x \in \widetilde{A}\) , let \(\gamma_{x}\) be the multiplication operator \(y \mapsto xy\) from A to A, and set \(\|x\|_{\widetilde{A}} = \| \gamma_{x}\|\) . The map \(x \mapsto \|x\|_{\widetilde{A}}\) is a semi-norm on \(\widetilde{A}\) . For all x and \(x'\) in \(\widetilde{A}\) , one has \(\|xx'\|_{\widetilde{A}} \leqslant \|x\|_{\widetilde{A}} \|x'\|_{\widetilde{A}}\) . Moreover, by Lemma 7, we have \(\|x\|_{\widetilde{A}} = \|x\|\) for all \(x \in \widetilde{A}\) .

Let us show that the map \(x \mapsto \|x\|_{\widetilde{\mathrm{A}}}\) is a norm on \(\widetilde{\mathrm{A}}\) . Let \(\lambda \in \mathbf{C}\) and \(x \in \mathrm{A}\) such that \(\| \lambda \widetilde{e} + x\|_{\widetilde{\mathrm{A}}} = 0\) . If \(\lambda \neq 0\) the condition \((\lambda \widetilde{e} + x)y = 0\) for all \(y \in \mathrm{A}\) implies that \(-\lambda^{-1}x\) is a left identity element in A. Likewise, the element \(-\overline{\lambda}^{-1}x^{*}\) is a right identity element. Thus, the algebra A would then have an identity element, contrary to the hypothesis. Hence, \(\lambda = 0\) . But then \(0 = \|x\|_{\widetilde{\mathrm{A}}} = \|x\|\) , and hence \(x = 0\) .

Since A is complete and of codimension 1 in \(\widetilde{\mathrm{A}}\) , the space \(\widetilde{\mathrm{A}}\) equipped with the norm \(x \mapsto \|x\|_{\widetilde{\mathrm{A}}}\) is complete. To conclude, it is therefore sufficient to show that one has \(\|x\|_{\widetilde{\mathrm{A}}}^{2} \leqslant \|x^{*}x\|_{\widetilde{\mathrm{A}}}\) for all \(x \in \widetilde{\mathrm{A}}\) (lemma 6). One may assume that \(\|x\|_{\widetilde{\mathrm{A}}} = 1\) . For every real number \(r < 1\) , there therefore exists \(y \in \mathrm{A}\) such that \(\|y\| = \|y^{*}\| \leqslant 1\) and \(\|xy\|^2 \geqslant r\) . Since \(xy \in \mathrm{A}\) , one has

\[
\| x ^ {*} x \| _ {\widetilde {A}} \geqslant \| x ^ {*} x y \| \geqslant \| y ^ {*} (x ^ {*} x) y \| = \| (x y) ^ {*} (x y) \| = \| x y \| ^ {2} \geqslant r.
\]

From this one deduces that \(\| x^{*}x\|_{\widetilde{\mathrm{A}}}\geqslant 1\) , and hence the involutive algebra \(\widetilde{\mathrm{A}}\) endowed with the norm \(x\mapsto \| x\|_{\widetilde{\mathrm{A}}}\) is a stellar algebra.

DEFINITION 4. --- One says that \(\widetilde{\mathrm{A}}\) , endowed with the norm of prop. 3, is the stellar algebra deduced from A by adjoining an identity element.

When \(A \neq \{0\}\) , the stellar-algebra norm on the normed involutive algebra \(\widetilde{A}\) is not the one considered in example 6 of I, p.~101 (cf.~exercise 10 of I, p.~181).

PROPOSITION 4. --- Let A be a stellar algebra.

\begin{enumerate}
\def\labelenumi{\alph{enumi})}
\item
  If A has a unit element and u is a unitary element of A, then \(\mathrm{Sp}(u) \subset \mathbf{U}\) ;
\item
  If h is a Hermitian element of A, then \(\mathrm{Sp}'(h) \subset \mathbf{R}\) .
\end{enumerate}

We may assume that A is nonzero. Let us prove the assertion \(a\) ). Let \(u\) be a unitary element of A. We have \(\|u\|=\|u^{-1}\|=1\) (remark 1), hence \(\mathrm{Sp}(u)\subset\mathbf{U}\) (I, p.~26, cor. 3). To prove \(b\) ), we may assume that A is unital (prop. 3). Let \(h\) be a Hermitian element of A. Then \(\exp(ih)\) is unitary (Lemma 5 of I, p.~102). Thus, by Cor. 2 of I, p.~67 and \(a\) ), we have \(\exp(i\mathrm{Sp}(h))=\mathrm{Sp}(\exp(ih))\subset\mathbf{U}\) , which means that \(\mathrm{Sp}(h)\subset\mathbf{R}\) .

PROPOSITION 5. — Let A be a unital stellar algebra and B a stellar subalgebra of A containing the identity element of A. Then B is a full subalgebra of A. In particular, one has \(\mathrm{Sp}_{\mathrm{B}}(x) = \mathrm{Sp}_{\mathrm{A}}(x)\) for all \(x\) in B.

Let x be a Hermitian element of B. Since \(\mathrm{Sp}_{\mathrm{B}}(x) \subset \mathbf{R}\) (prop. 4), prop. 6 of I, p.~28 shows that \(\mathrm{Sp}_{\mathrm{B}}(x) = \mathrm{Sp}_{\mathrm{A}}(x)\) . In particular, x is invertible in B if and only if it is invertible in A.

Now let x be any element of B that is invertible in A. Then \(x^{*}\) is invertible in A, and \(xx^{*}\) is invertible in A. Since \(xx^{*}\) is Hermitian, the preceding shows that \(xx^{*}\) is invertible in B. This implies that x is right-invertible in B. Similarly, one verifies that x is left-invertible in B, and consequently that x is invertible in B. Thus, B is a full subalgebra of A.

COROLLARY. --- Let A be a stellar algebra and let B be a stellar subalgebra of A. Then we have \(\mathrm{Sp}_{\mathrm{B}}^{\prime}(x)=\mathrm{Sp}_{\mathrm{A}}^{\prime}(x)\) for all x in B.

By adjoining an identity element (prop. 3), this follows from prop. 5.

PROPOSITION 6 (Fuglede–Putnam Theorem)

Let A be a unital stellar algebra. Let a and b be normal elements of A. If \(c \in A\) satisfies ac = cb, then \(a^{*}c = cb^{*}\) .

The hypothesis implies \((wa)^{k}c = c(wb)^{k}\) for every integer \(k \geqslant 0\) and every \(w \in C\) , hence \(e^{wa}c = ce^{wb}\) (formula (9) of I, p.~78). Consider the function f from C to A defined by \(z \mapsto e^{-za^{*}}ce^{zb^{*}}\) . This is a holomorphic function on C, whose derivative satisfies

\[
f ^ {\prime} (z) = - a ^ {*} e ^ {- z a ^ {*}} c e ^ {z b ^ {*}} + e ^ {- z a ^ {*}} c b ^ {*} e ^ {z b ^ {*}}
\]

for all \(z \in \mathbf{C}\) . Since \(c = e^{-\overline{z}a} ce^{\overline{z}b}\) one can write

\[
f (z) = e ^ {- z a ^ {*}} e ^ {\overline {{z}} a} c e ^ {- \overline {{z}} b} e ^ {z b ^ {*}}.
\]

Since \(a\) and \(b\) are normal, the elements \(e^{-za^*}e^{\overline{z} a} = e^{-za^* + \overline{z} a}\) and \(e^{-\overline{z} b}e^{zb^*} = e^{-\overline{z} b + zb^*}\) of A are unitary for all \(z\in \mathbf{C}\) (Lemma 5 of I, p.~102), hence of norm 1. Consequently, one has \(\| f(z)\| \leqslant \| c\|\) for all \(z\in \mathbf{C}\) . The function \(f\) is therefore constant (VAR, R1, 3.3.6, p.~29), that is, \(f(z) = f(0) = c\) for all \(z\in \mathbf{C}\) . But then \(-a^{*}c + cb^{*} = f'(0) = 0\) .

COROLLARY. --- Let A be a stellar algebra and let a be a normal element of A. The commutant (resp. bicommutant) of \(\{a, a^{*}\}\) coincides with the commutant (resp. bicommutant) of a.

It suffices to take b = a in the proposition.

